\section{Operator Precedence Grammar}
\label{sec:ch11-operator}

\subsection{Definisi}

Operator precedence grammar adalah grammar yang memenuhi:
\begin{enumerate}[label=\arabic*.,leftmargin=*]
  \item tidak memiliki produksi \(\varepsilon\);
  \item tidak memiliki dua nonterminal bersebelahan pada ruas kanan.
\end{enumerate}

Sifat ini memungkinkan parser menentukan hubungan precedence antara
terminal yang berdekatan tanpa membangun tabel LR lengkap.

\subsection{Relasi Precedence}

Tiga relasi khusus digunakan:
\begin{itemize}[leftmargin=*]
  \item \(a <\cdot b\): \(a\) memiliki precedence lebih rendah
        daripada \(b\);
  \item \(a \doteq b\): \(a\) dan \(b\) berada pada tingkat yang
        sama;
  \item \(a \cdot> b\): \(a\) memiliki precedence lebih tinggi
        daripada \(b\).
\end{itemize}
Simbol \(<\), \(=\), dan \(>\) di sini bukan operator aritmetika,
melainkan relasi precedence.

\subsection{Tabel Operator Precedence}

Untuk grammar ekspresi dengan terminal
\(\{\code{id}, +, *, \$, (, )\}\), tabelnya berbentuk:
\begin{table}[H]
\centering
\caption{Tabel operator precedence sederhana}
\label{tab:ch11-precedence}
\begin{tabular}{@{}c c c c c c@{}}
\toprule
 & \code{id} & \code{+} & \code{*} & \code{(} & \code{)} \\
\midrule
\code{id} & & \(\cdot>\) & \(\cdot>\) & & \(\cdot>\) \\
\code{+} & \(<\cdot\) & \(\cdot>\) & \(<\cdot\) & \(<\cdot\) & \(\cdot>\) \\
\code{*} & \(<\cdot\) & \(\cdot>\) & \(\cdot>\) & \(<\cdot\) & \(\cdot>\) \\
\code{(} & \(<\cdot\) & \(<\cdot\) & \(<\cdot\) & \(<\cdot\) & \(\doteq\) \\
\code{)} & & \(\cdot>\) & \(\cdot>\) & & \(\cdot>\) \\
\bottomrule
\end{tabular}
\end{table}

Sel kosong menandakan kesalahan sintaksis. Tabel ini lebih kecil
daripada tabel LR karena hanya memuat terminal, bukan himpunan item.

\subsection{Contoh Relasi}

Untuk ekspresi \(\code{id} + \code{id} * \code{id}\):
\begin{itemize}[leftmargin=*]
  \item \(\code{id} \cdot> +\): \code{id} harus direduksi sebelum
        \(+\);
  \item \(+ <\cdot \code{id}\): \(+\) menunggu operand;
  \item \(+ <\cdot *\): \(*\) memiliki precedence lebih tinggi;
  \item \(* \cdot> \code{id}\): \(*\) menunggu operand.
\end{itemize}

Relasi ini menentukan kapan parser harus shift dan kapan harus
reduce. Parser membandingkan terminal teratas stack dengan token
saat ini.

\subsection{Keterbatasan}

Operator precedence hanya menangani bahasa dengan struktur operator
yang sederhana. Grammar dengan konstruksi seperti dangling else,
deklarasi, atau pemanggilan fungsi memerlukan parser LR yang lebih
umum. Karena itu, operator precedence lebih cocok untuk ekspresi
daripada seluruh bahasa pemrograman.
